% Options for packages loaded elsewhere
\PassOptionsToPackage{unicode}{hyperref}
\PassOptionsToPackage{hyphens}{url}
\documentclass[
]{article}
\usepackage{xcolor}
\usepackage{amsmath,amssymb}
\setcounter{secnumdepth}{-\maxdimen} % remove section numbering
\usepackage{iftex}
\ifPDFTeX
  \usepackage[T1]{fontenc}
  \usepackage[utf8]{inputenc}
  \usepackage{textcomp} % provide euro and other symbols
\else % if luatex or xetex
  \usepackage{unicode-math} % this also loads fontspec
  \defaultfontfeatures{Scale=MatchLowercase}
  \defaultfontfeatures[\rmfamily]{Ligatures=TeX,Scale=1}
\fi
\usepackage{lmodern}
\ifPDFTeX\else
  % xetex/luatex font selection
\fi
% Use upquote if available, for straight quotes in verbatim environments
\IfFileExists{upquote.sty}{\usepackage{upquote}}{}
\IfFileExists{microtype.sty}{% use microtype if available
  \usepackage[]{microtype}
  \UseMicrotypeSet[protrusion]{basicmath} % disable protrusion for tt fonts
}{}
\makeatletter
\@ifundefined{KOMAClassName}{% if non-KOMA class
  \IfFileExists{parskip.sty}{%
    \usepackage{parskip}
  }{% else
    \setlength{\parindent}{0pt}
    \setlength{\parskip}{6pt plus 2pt minus 1pt}}
}{% if KOMA class
  \KOMAoptions{parskip=half}}
\makeatother
\usepackage{color}
\usepackage{fancyvrb}
\newcommand{\VerbBar}{|}
\newcommand{\VERB}{\Verb[commandchars=\\\{\}]}
\DefineVerbatimEnvironment{Highlighting}{Verbatim}{commandchars=\\\{\}}
% Add ',fontsize=\small' for more characters per line
\newenvironment{Shaded}{}{}
\newcommand{\AlertTok}[1]{\textcolor[rgb]{1.00,0.00,0.00}{\textbf{#1}}}
\newcommand{\AnnotationTok}[1]{\textcolor[rgb]{0.38,0.63,0.69}{\textbf{\textit{#1}}}}
\newcommand{\AttributeTok}[1]{\textcolor[rgb]{0.49,0.56,0.16}{#1}}
\newcommand{\BaseNTok}[1]{\textcolor[rgb]{0.25,0.63,0.44}{#1}}
\newcommand{\BuiltInTok}[1]{\textcolor[rgb]{0.00,0.50,0.00}{#1}}
\newcommand{\CharTok}[1]{\textcolor[rgb]{0.25,0.44,0.63}{#1}}
\newcommand{\CommentTok}[1]{\textcolor[rgb]{0.38,0.63,0.69}{\textit{#1}}}
\newcommand{\CommentVarTok}[1]{\textcolor[rgb]{0.38,0.63,0.69}{\textbf{\textit{#1}}}}
\newcommand{\ConstantTok}[1]{\textcolor[rgb]{0.53,0.00,0.00}{#1}}
\newcommand{\ControlFlowTok}[1]{\textcolor[rgb]{0.00,0.44,0.13}{\textbf{#1}}}
\newcommand{\DataTypeTok}[1]{\textcolor[rgb]{0.56,0.13,0.00}{#1}}
\newcommand{\DecValTok}[1]{\textcolor[rgb]{0.25,0.63,0.44}{#1}}
\newcommand{\DocumentationTok}[1]{\textcolor[rgb]{0.73,0.13,0.13}{\textit{#1}}}
\newcommand{\ErrorTok}[1]{\textcolor[rgb]{1.00,0.00,0.00}{\textbf{#1}}}
\newcommand{\ExtensionTok}[1]{#1}
\newcommand{\FloatTok}[1]{\textcolor[rgb]{0.25,0.63,0.44}{#1}}
\newcommand{\FunctionTok}[1]{\textcolor[rgb]{0.02,0.16,0.49}{#1}}
\newcommand{\ImportTok}[1]{\textcolor[rgb]{0.00,0.50,0.00}{\textbf{#1}}}
\newcommand{\InformationTok}[1]{\textcolor[rgb]{0.38,0.63,0.69}{\textbf{\textit{#1}}}}
\newcommand{\KeywordTok}[1]{\textcolor[rgb]{0.00,0.44,0.13}{\textbf{#1}}}
\newcommand{\NormalTok}[1]{#1}
\newcommand{\OperatorTok}[1]{\textcolor[rgb]{0.40,0.40,0.40}{#1}}
\newcommand{\OtherTok}[1]{\textcolor[rgb]{0.00,0.44,0.13}{#1}}
\newcommand{\PreprocessorTok}[1]{\textcolor[rgb]{0.74,0.48,0.00}{#1}}
\newcommand{\RegionMarkerTok}[1]{#1}
\newcommand{\SpecialCharTok}[1]{\textcolor[rgb]{0.25,0.44,0.63}{#1}}
\newcommand{\SpecialStringTok}[1]{\textcolor[rgb]{0.73,0.40,0.53}{#1}}
\newcommand{\StringTok}[1]{\textcolor[rgb]{0.25,0.44,0.63}{#1}}
\newcommand{\VariableTok}[1]{\textcolor[rgb]{0.10,0.09,0.49}{#1}}
\newcommand{\VerbatimStringTok}[1]{\textcolor[rgb]{0.25,0.44,0.63}{#1}}
\newcommand{\WarningTok}[1]{\textcolor[rgb]{0.38,0.63,0.69}{\textbf{\textit{#1}}}}
\usepackage{longtable,booktabs,array}
\newcounter{none} % for unnumbered tables
\usepackage{calc} % for calculating minipage widths
% Correct order of tables after \paragraph or \subparagraph
\usepackage{etoolbox}
\makeatletter
\patchcmd\longtable{\par}{\if@noskipsec\mbox{}\fi\par}{}{}
\makeatother
% Allow footnotes in longtable head/foot
\IfFileExists{footnotehyper.sty}{\usepackage{footnotehyper}}{\usepackage{footnote}}
\makesavenoteenv{longtable}
\setlength{\emergencystretch}{3em} % prevent overfull lines
\providecommand{\tightlist}{%
  \setlength{\itemsep}{0pt}\setlength{\parskip}{0pt}}
\usepackage{bookmark}
\IfFileExists{xurl.sty}{\usepackage{xurl}}{} % add URL line breaks if available
\urlstyle{same}
\hypersetup{
  hidelinks,
  pdfcreator={LaTeX via pandoc}}

\author{}
\date{}

\begin{document}

\section{Whose Vote Should Count More: Optimal Integration of Labels
from Labelers of Unknown
Expertise}\label{whose-vote-should-count-more-optimal-integration-of-labels-from-labelers-of-unknown-expertise}

Jacob Whitehill, Paul Ruvolo, Tingfan Wu, Jacob Bergsma, and Javier
Movellan

Machine Perception Laboratory

University of California, San Diego

La Jolla, CA, USA

\{ jake, paul, ting, jbergsma, movellan \}@mplab.ucsd.edu

\section{Abstract}\label{abstract}

Modern machine learning-based approaches to computer vision require very
large databases of hand labeled images. Some contemporary vision systems
already require on the order of millions of images for training (e.g.,
Omron face detector {[}9{]}). New Internet-based services allow for a
large number of labelers to collaborate around the world at very low
cost. However, using these services brings interesting theoretical and
practical challenges: (1) The labelers may have wide ranging levels of
expertise which are unknown a priori, and in some cases may be
adversarial; (2) images may vary in their level of difficulty; and (3)
multiple labels for the same image must be combined to provide an
estimate of the actual label of the image. Probabilistic approaches
provide a principled way to approach these problems. In this paper we
present a probabilistic model and use it to simultaneously infer the
label of each image, the expertise of each labeler, and the difficulty
of each image. On both simulated and real data, we demonstrate that the
model outperforms the commonly used ``Majority Vote'' heuristic for
inferring image labels, and is robust to both noisy and adversarial
labelers.

\section{1 Introduction}\label{introduction}

In recent years machine learning-based approaches to computer vision
have helped to greatly accelerate progress in the field. However, it is
now becoming clear that many practical applications require very large
databases of hand labeled images. The labeling of very large datasets is
becoming a bottleneck for progress. One approach to address this
incoming problem is to make use of the vast human resources on the
Internet. Indeed, projects like the ESP game {[}17{]}, the Listen
game{[}16{]}, Soylent Grid {[}15{]}, and reCAPTCHA {[}18{]} have
revealed the possibility of harnessing human resources to solve
difficult machine learning problems. While these approaches use clever
schemes to obtain data from humans for free, a more direct approach is
to hire labelers online. Recent Web tools such as Amazon's Mechanical
Turk {[}1{]} provide ideal solutions for high-speed, low cost labeling
of massive databases.

Due to the distributed and anonymous nature of these tools, interesting
theoretical and practical challenges arise. For example, principled
methods are needed to combine the labels from multiple experts and to
estimate the certainty of the current labels. Which image should be
labeled (or relabeled) next must also be decided -- it may be prudent,
for example, to collect many labels for each image in order to increase
one's confidence in that image's label. However, if an image is easy and
the labelers of that image are reliable, a few labels may be sufficient
and valuable resources may be used to label other images. In practice,
combining the labels of multiple coders is a challenging process due to
the fact that: (1) The labelers may have wide ranging levels of
expertise which are

unknown a priori, and in some cases may be adversarial; (2) images may
also vary in their level of difficulty, in a manner that may also be
unknown a priori.

Probabilistic methods provide a principled way to approach this problem
using standard inference tools. We explore one such approach by
formulating a probabilistic model of the labeling process, which we call
GLAD (Generative model of Labels, Abilities, and Difficulties), and
using inference methods to simultaneously infer the expertise of each
labeler, the difficulty of each image, and the most probable label for
each image. On both simulated and real-life data, we demonstrate that
the model outperforms the commonly used ``Majority Vote'' heuristic for
inferring image labels, and is robust to both adversarial and noisy
labelers.

\section{2 Modeling the Labeling
Process}\label{modeling-the-labeling-process}

Consider a database of \(n\) images, each of which belongs to one of two
possible categories of interest (e.g., face/non-face; male/female;
smile/non-smile; etc.). We wish to determine the class label \(Z_{j}\)
(0 or 1) of each image \(j\) by querying from \(m\) labelers. The
observed labels depend on several causal factors: (1) the difficulty of
the image; (2) the expertise of the labeler; and (3) the true label. We
model the difficulty of image \(j\) using the parameter
\(1 / \beta_{j} \in [0, \infty)\) where \(\beta_{j}\) is constrained to
be positive. Here \(1 / \beta_{j} = \infty\) means the image is very
ambiguous and hence even the most proficient labeler has a \(50\%\)
chance of labeling it correctly. \(1 / \beta_{j} = 0\) means the image
is so easy that even the most obtuse labeler will always label it
correctly.

The expertise of each labeler i is modeled by the parameter
\(\alpha_{i} \in (-\infty, +\infty)\) . Here an \(\alpha = +\infty\)
means the labeler always labels images correctly; \(-\infty\) means the
labeler always labels the images incorrectly, i.e., he/she can
distinguish between the two classes perfectly but always inverts the
label, either maliciously or because of a consistent misunderstanding.
In this case ( \(\alpha_{i} < 0\) ), the labeler is said to be
adversarial. Finally, \(\alpha_{i} = 0\) means that the labeler cannot
discriminate between the two classes -- his/her labels carry no
information about the true image label \(Z_{j}\) . Note that we do not
require the labelers to be human -- labelers can also be, for instance,
automatic classifiers. Hence, the proposed approach will provide a
principled way of combining labels from any combination of human and
previously existing machine-based classifiers.

The labels given by labeler i to image j (which we call the given
labels) are denoted as \(L_{ij}\) and, under the model, are generated as
follows:

\[
p \left(L _ {i j} = Z _ {j} \mid \alpha_ {i}, \beta_ {j}\right) = \frac {1}{1 + e ^ {- \alpha_ {i} \beta_ {j}}} \tag {1}
\]

Thus, under the model, the log odds for the obtained labels being
correct are a bilinear function function of the difficulty of the label
and the expertise of the labeler, i.e.,

\[
\log \frac {p \left(L _ {i j} = Z _ {j}\right)}{1 - p \left(L _ {i j} = Z _ {j}\right)} = \alpha_ {i} \beta_ {j} \tag {2}
\]

More skilled labelers (higher \(\alpha_{i}\) ) have a higher probability
of labeling correctly. As the difficulty \(1/\beta_{j}\) of an image
increases, the probability of the label being correct moves toward 0.5.
Similarly, as the labeler's expertise decreases (lower \(\alpha_{i}\) ),
the chance of correctness likewise drops to 0.5. Adversarial labelers
are simply labelers with negative \(\alpha\) .

Figure 1 shows the causal structure of the model. True image labels
\(Z_{j}\) , labeler accuracy values \(\alpha_{i}\) , and image
difficulty values \(\beta_{j}\) are sampled from a known prior
distribution. These determine the observed labels according to Equation
1. Given a set of observed labels \(l = \{l_{ij}\}\) , the task is to
infer simultaneously the most likely values of \(Z = \{Z_{j}\}\) (the
true image labels) as well as the labeler accuracies
\(\alpha = \{\alpha_{i}\}\) and the image difficulty parameters
\(\beta = \{\beta_{j}\}\) . In the next section we derive the Maximum
Likelihood algorithm for inferring these values.

\section{3 Inference}\label{inference}

The observed labels are samples from the \(\{L_{ij}\}\) random
variables. The unobserved variables are the true image labels \(Z_{j}\)
, the different labeler accuracies \(\alpha_{i}\) , and the image
difficulty parameters \(1/\beta_{j}\) . Our goal is to efficiently
search for the most probable values of the unobservable variables

\textit{[Image omitted: images/0572f04708cdde6404bc065c52d874b4e95a44935af47972db8249bd280646c3.jpg]}

flowchart

\begin{Shaded}
\begin{Highlighting}[]
\NormalTok{graph TD}
\NormalTok{    subgraph Image difficulties}
\NormalTok{        A1["β₁"] {-}{-}\textgreater{} B1["z₁"]}
\NormalTok{        A1 {-}{-}\textgreater{} C1["L₁₁"]}
\NormalTok{        A1 {-}{-}\textgreater{} D1["L₂₁"]}
\NormalTok{        A1 {-}{-}\textgreater{} E1["..."]}
\NormalTok{        A1 {-}{-}\textgreater{} F1["L₁₂"]}
\NormalTok{        A1 {-}{-}\textgreater{} G1["L₂₂"]}
\NormalTok{        A1 {-}{-}\textgreater{} H1["L₃₂"]}
\NormalTok{        A2["β₂"] {-}{-}\textgreater{} B2["z₂"]}
\NormalTok{        A2 {-}{-}\textgreater{} C2["L₂₁"]}
\NormalTok{        A2 {-}{-}\textgreater{} D2["L₁₂"]}
\NormalTok{        A2 {-}{-}\textgreater{} E2["L₂₂"]}
\NormalTok{        A2 {-}{-}\textgreater{} F2["L₃₂"]}
\NormalTok{        A2 {-}{-}\textgreater{} G2["L₃₂"]}
\NormalTok{        A2 {-}{-}\textgreater{} H2["L₃₂"]}
\NormalTok{        A2 {-}{-}\textgreater{} I2["..."]}
\NormalTok{        A2 {-}{-}\textgreater{} J2["L₃₂"]}
\NormalTok{    end}

\NormalTok{    subgraph True labels}
\NormalTok{        K1["z₃"] {-}{-}\textgreater{} L1["z₃"]}
\NormalTok{        K1 {-}{-}\textgreater{} M1["..."]}
\NormalTok{        K2["zₙ"] {-}{-}\textgreater{} L2["zₙ"]}
\NormalTok{    end}

\NormalTok{    subgraph Observed labels}
\NormalTok{        N1["..."]}
\NormalTok{        N2["..."]}
\NormalTok{        N3["..."]}
\NormalTok{        N4["..."]}
\NormalTok{        N5["..."]}
\NormalTok{        N6["..."]}
\NormalTok{        N7["..."]}
\NormalTok{        N8["..."]}
\NormalTok{        N9["..."]}
\NormalTok{        N10["α₁"]}
\NormalTok{        N11["α₂"]}
\NormalTok{        N12["α₃"]}
\NormalTok{    end}

\NormalTok{    subgraph Labeler accuracies}
\NormalTok{        O1["..."]}
\NormalTok{        O2["αₘ"]}
\NormalTok{    end}
\end{Highlighting}
\end{Shaded}

Figure 1: Graphical model of image difficulties, true image labels,
observed labels, and labeler accuracies. Only the shaded variables are
observed.

Z, \(\alpha\) and \(\beta\) given the observed data. Here we can use
Expectation-Maximization approach (EM) to obtain maximum likelihood
estimates of the parameters of interest (the full derivation is in the
Supplementary Materials):

E step: Let the set of all given labels for an image j be denoted as
\(l_{j} = \{l_{ij'} \mid j' = j\}\) . Note that not every labeler must
label every single image. In this case, the index variable i in
\(l_{ij'}\) refers only to those labelers who labeled image j. We need
to compute the posterior probabilities of all \(z_{j} \in \{0, 1\}\)
given the \(\alpha, \beta\) values from the last M step and the observed
labels:

\[
\begin{array}{l} p (z _ {j} | \mathbf {l}, \boldsymbol {\alpha}, \boldsymbol {\beta}) = p (z _ {j} | \mathbf {l} _ {j}, \boldsymbol {\alpha}, \boldsymbol {\beta} _ {j}) \\ \propto \quad p (z _ {j} | \boldsymbol {\alpha}, \beta_ {j}) p (\mathbf {l} _ {j} | z _ {j}, \boldsymbol {\alpha}, \beta_ {j}) \\ \propto \quad p (z _ {j}) \prod_ {i} p (l _ {i j} | z _ {j}, \alpha_ {i}, \beta_ {j}) \\ \end{array}
\]

where we noted that \(p(z_{j}|\alpha, \beta_{j}) = p(z_{j})\) using the
conditional independence assumptions from the graphical model.

M step: We maximize the standard auxiliary function Q, which is defined
as the expectation of the joint log-likelihood of the observed and
hidden variables \((1, \mathbf{Z})\) given the parameters
\((\alpha, \beta)\) , w.r.t. the posterior probabilities of the Z values
computed during the last E step:

\[
\begin{array}{l} Q (\boldsymbol {\alpha}, \boldsymbol {\beta}) = E [ \ln p (\mathbf {l}, \mathbf {z} | \boldsymbol {\alpha}, \boldsymbol {\beta}) ] \\ = E \left[ \ln \prod_ {j} \left(p \left(z _ {j}\right) \prod_ {i} p \left(l _ {i j} \mid z _ {j}, \alpha_ {i}, \beta_ {j}\right)\right) \right] \\ = \sum_ {j} E [ \ln p (z _ {j}) ] + \sum_ {i j} E [ \ln p (l _ {i j} | z _ {j}, \alpha_ {i}, \beta_ {j}) ] \\ \end{array}
\]

since \(l_{ij}\) are cond. indep. given \(z, \alpha, \beta\)

where the expectation is taken over z given the old parameter values
\(\alpha^{old}, \beta^{old}\) as estimated during the last E-step. Using
gradient ascent, we find values of \(\alpha\) and \(\beta\) that locally
maximize Q.

\section{\texorpdfstring{3.1 Priors on
\(\alpha, \beta\)}{3.1 Priors on \textbackslash alpha, \textbackslash beta}}\label{priors-on-alpha-beta}

The Q function can be modified straightforwardly to handle a prior over
each \(\alpha_{i}\) and \(\beta_{j}\) by adding a log-prior term for
each of these variables. These priors may be useful, for example, if we
know that most labelers are not adversarial. In this case, the prior for
\(\alpha\) can be made very low for \(\alpha < 0\) .

The prior probabilities are also useful when the ground-truth Z value of
particular images is (some-how) known for certain. By ``clamping'' the Z
values (using the prior) for the images on which the

true label is known for sure, the model may be able to better estimate
the other parameters. The Z values for such images can be clamped by
setting the prior probability \(p(z_{j})\) (used in the E-Step) for
these images to be very high towards one particular class. In our
implementation we used Gaussian priors ( \(\mu = 1, \sigma = 1\) ) for
\(\alpha\) . For \(\beta\) , we need a prior that does not generate
negative values. To do so we re-parameterized
\(\beta \doteq e^{\beta'}\) and imposed a Gaussian prior (
\(\mu = 1, \sigma = 1\) ) on \(\beta'\) .

\section{3.2 Computational Complexity}\label{computational-complexity}

The computational complexity of the E-Step is linear in the number of
images and the total number of labels. For the M-Step, the values of Q
and \(\nabla Q\) must be computed repeatedly until convergence. \(^{1}\)
Computing each function is linear in the number of images, number of
labelers, and total number of image labels.

Empirically when using the approach on a database of 1 million images
that we recently collected and labeled we found that the EM procedure
converged in about 10 minutes using a single core of a Xeon 2.8 GHz
processor. The algorithm is parallelizable and hence this running time
could be reduced substantially using multiple cores. Real time inference
may also be possible if we maintain parameters close to the solution
that are updated as new labels become available. This would allow using
the algorithm in an active manner to choose in real-time which images
should be labeled next so as to minimize the uncertainty about the image
labels.

\section{4 Simulations}\label{simulations}

Here we explore the performance of the model using a set of image labels
generated by the model itself. Since, in this case we know the
parameters Z, \(\alpha\) , and \(\beta\) that generated the observed
labels, we can compare them with corresponding parameters estimated
using the EM procedure.

In particular, we simulated between 4 and 20 labelers, each labeling
2000 images, whose true labels Z were either 0 or 1 with equal
probability. The accuracy \(\alpha_{i}\) of each labeler was drawn from
a normal distribution with mean 1 and variance 1. The inverse-difficulty
for each image \(\beta_{j}\) was generated by exponentiating a draw from
a normal distribution with mean 1 and variance 1. Given these labeler
abilities and image difficulties, the observed labels \(l_{ij}\) were
sampled according to Equation 1 using Z. Finally, the EM inference
procedure described above was executed to estimate \(\alpha\) ,
\(\beta\) , Z. This procedure was repeated 40 times to smooth out
variability between trials. On each trial we computed the correlation
between the parameter estimates \(\hat{\alpha}\) , \(\hat{\beta}\) and
the true parameter values \(\alpha\) , \(\beta\) . The results (averaged
over all 40 experimental runs) are shown in Figure 2. As expected, as
the number of labelers grows, the parameter estimates converge to the
true values.

We also computed the proportion of label estimates \(\hat{Z}\) that
matched the true image labels Z. We compared the maximum likelihood
estimates of the GLAD model to estimates obtained by taking the majority
vote as the predicted label. The predictions of the proposed GLAD model
were obtained by thresholding at 0.5 the posterior probability of the
label of each image being of class 1 given the accuracy and difficulty
parameters returned by EM (see Section 3). Results are shown in Figure
2. GLAD makes fewer errors than the majority vote heuristic. The
difference between the two approaches is particularly pronounced when
the number of labelers per image is small. On many images, GLAD
correctly infers the true image label Z even when that Z value was the
minority opinion. In essence, GLAD is exploiting the fact that some
labelers are experts (which it infers automatically), and hence their
votes should count more on these images than the votes of less skilled
labelers.

Modeling Image Difficulty : To explore the importance of estimating
image difficulty we performed a simple simulation: Image labels (0 or 1)
were assigned randomly (with equal probability) to 1000 images. Half of
the images were ``hard'', and half were ``easy.'' Fifty simulated
labelers labeled all 1000 images. The proportion of ``good'' to ``bad''
labelers is 25:1. The probability of correctness for each image
difficulty and labeler quality combination was given by the table below:

\textit{[Image omitted: images/bd360b3404c8c867c571bc1f32c4298135e3f3c66825d0420d9cb20a9292370e.jpg]}

line

{\def\LTcaptype{none} % do not increment counter
\begin{longtable}[]{@{}lll@{}}
\toprule\noalign{}
Number of Labelers & GLAD & Majority vote \\
\midrule\noalign{}
\endhead
\bottomrule\noalign{}
\endlastfoot
5 & 0.90 & 0.85 \\
10 & 0.97 & 0.93 \\
15 & 0.98 & 0.96 \\
20 & 0.99 & 0.97 \\
\end{longtable}
}

\textit{[Image omitted: images/37e28f98e5c24567eb74777003830b187818516ab0ac3fc138cd8dd68170ebec.jpg]}

line

{\def\LTcaptype{none} % do not increment counter
\begin{longtable}[]{@{}lll@{}}
\toprule\noalign{}
Number of Labelers & Beta: Spearman Corr. & Alpha: Pearson Corr. \\
\midrule\noalign{}
\endhead
\bottomrule\noalign{}
\endlastfoot
5 & 0.5 & 0.85 \\
10 & 0.75 & 0.95 \\
15 & 0.85 & 0.98 \\
20 & 0.88 & 0.99 \\
\end{longtable}
}

Figure 2: Left: The accuracies of the GLAD model versus simple voting
for inferring the underlying class labels on simulation data. Right: The
ability of GLAD to recover the true alpha and beta parameters on
simulation data.

Labeler type

Image Type

Hard

Easy

Good

0.95

1

Bad

0.54

1

We measured performance in terms of proportion of correctly estimated
labels. We compared three approaches: (1) our proposed method, GLAD; (2)
the method proposed in {[}5{]}, which models labeler ability but not
image difficulty; and (3) Majority Vote. The simulations were repeated
20 times and average performance calculated for the three methods. The
results shown below indicated that modeling image difficulty can result
in significant performance improvements.

Method

Error

GLAD

4.5\%

Majority Vote

11.2\%

Dawid \& Skene {[}5{]}

8.4\%

\section{4.1 Stability of EM under Various Starting
Points}\label{stability-of-em-under-various-starting-points}

Empirically we found that the EM procedure was fairly insensitive to
varying the starting point of the parameter values. In a simulation
study of 2000 images and 20 labelers, we randomly selected each
\(\alpha_{i} \sim U[0,4]\) and \(\log(\beta_{j}) \sim U[0,3]\) , and EM
was run until convergence. Over the 50 simulation runs, the average
percent-correct of the inferred labels was 85.74\%, and the standard
deviation of the percent-correct over all the trials was only 0.024\%.

\section{5 Empirical Study I:
Greebles}\label{empirical-study-i-greebles}

As a first test-bed for GLAD using real data obtained from the
Mechanical Turk, we posted pictures of 100 ``Greebles'' {[}6{]}, which
are synthetically generated images that were originally created to study
human perceptual expertise. Greebles somewhat resemble human faces and
have a ``gender'': Males have horn-like organs that point up, whereas
for females the horns point down. See Figure 3 (left) for examples. Each
of the 100 Greeble images was labeled by 10 different human coders on
the Turk for gender (male/female). Four greebles of each gender
(separate from the 100 labeled images) were given as examples of each
class. Shown at a resolution of 48x48 pixels, the task required careful
inspection of the images in order to label them correctly. The
ground-truth gender values were all known with certainty (since they are
rendered objects) and thus provided a means of measuring the accuracy of
inferred image labels.

\textit{[Image omitted: images/24c136103c89ea3df71638fe4c2fc7bda6062fd0d994ca8f1d296a934aeb44bd.jpg]}

line

{\def\LTcaptype{none} % do not increment counter
\begin{longtable}[]{@{}lll@{}}
\toprule\noalign{}
Number of labels per image & GLAD & Majority Vote \\
\midrule\noalign{}
\endhead
\bottomrule\noalign{}
\endlastfoot
2 & 0.95 & 0.85 \\
3 & 0.97 & 0.95 \\
4 & 0.98 & 0.93 \\
5 & 0.99 & 0.97 \\
6 & 0.995 & 0.96 \\
7 & 0.998 & 0.98 \\
8 & 1.0 & 0.98 \\
\end{longtable}
}

Figure 3: Left: Examples of Greebles. The top two are ``male'' and the
bottom two are ``female.'' Right: Accuracy of the inferred labels, as a
function of the number of labels M obtained for each image, of the
Greeble images using either GLAD or Majority Vote. Results were averaged
over 100 experimental runs.

We studied the effect of varying the number of labels M obtained from
different labelers for each image, on the accuracy of the inferred Z.
Hence, from the 10 labels total we obtained per Greeble image, we
randomly sampled \(2 \leq M \leq 8\) labels over all labelers during
each experimental trial. On each trial we compared the accuracy of
labels Z as estimated by GLAD (using a threshold of 0.5 on \(p(Z)\) ) to
labels as estimated by the Majority Vote heuristic. For each value of M
we averaged performance for each method over 100 trials.

Results are shown in Figure 3 (right). For all values of M we tested,
the labels as inferred by GLAD are significantly higher than for
Majority Vote (p \textless{} 0.01). This means that, in order to achieve
the same level of accuracy, fewer labels are needed. Moreover, the
variance in accuracy was less for GLAD than for Majority Vote for all M
that were tested, suggesting that the quality of GLAD's outputs is more
stable than of the heuristic method. Finally, notice how, for the even
values of M, the Majority Vote accuracy decreases. This may stem from
the lack of optimal decision rule under Majority Vote when an equal
number of labelers say an image is Male as who say it is Female. GLAD,
since it makes its decisions by also taking ability and difficulty into
account, does not suffer from this problem.

\section{6 Empirical Study II: Duchenne
Smiles}\label{empirical-study-ii-duchenne-smiles}

As a second experiment, we used the Mechanical Turk to label face images
containing smiles as either Duchenne or Non-Duchenne. A Duchenne smile
(``enjoyment'' smile) is distinguished from a Non-Duchenne (``social''
smile) through the activation of the Orbicularis Oculi muscle around the
eyes, which the former exhibits and the latter does not (see Figure 4
for examples). Distinguishing the two kinds of smiles has applications
in various domains including psychology experiments, human-computer
interaction, and marketing research. Reliable coding of Duchenne smiles
is a difficult task even for certified experts in the Facial Action
Coding System.

We obtained Duchenne/Non-Duchenne labels for 160 images from 20
different Mechanical Turk labelers; in total, there were 3572 labels.
(Hence, labelers labeled each image a variable number of times.) For
ground truth, these images were also labeled by two certified experts in
the Facial Action Coding System. According to the expert labels, 58 out
of 160 images contained Duchenne smiles.

Using the labels obtained from the Mechanical Turk, we inferred the
image labels using either GLAD or the Majority Vote heuristic, and then
compared them to ground truth.

Duchenne Smiles\\
\textit{[Image omitted: images/bf64d411828c370c5ca65d3775e790678c97d806913418988c461d6ecfdf9d5f.jpg]}

\textit{[Image omitted: images/5e975e659aa816053309224f1386e9404e6748a5bf7d9974f60ea0575f697e52.jpg]}

\textit{[Image omitted: images/2d6a09eecc579c648d376d3abf8264b3c5800ee1f8f35f4eba8ec1ced5e32a3b.jpg]}

Non-Duchenne Smiles\\
\textit{[Image omitted: images/442b9291cbabaf91587564a303ba27ee3bca989f73a017eaafab86178fa379c9.jpg]}

\textit{[Image omitted: images/2e7257e6de1aa83bcdcabae92668aa85a560a8a056cc9cb3167acfd254bb9f5c.jpg]}

\textit{[Image omitted: images/3271fbe3c7d014415133f14ca3e8602d0973f967d2d4679900866d7e2b813370.jpg]}\\
Figure 4: Examples of Duchenne (left) and Non-Duchenne (right) smiles.
The distinction lies in the activation of Orbicularis Oculi muscle
around the eyes, and is difficult to discriminate even for experts.

\textit{[Image omitted: images/4f166b36afc7c91d92b2f89ae6c63eb81588794e1280acd6a17886f232234a36.jpg]}

line

{\def\LTcaptype{none} % do not increment counter
\begin{longtable}[]{@{}lll@{}}
\toprule\noalign{}
Num of Noisy Labels & GLAD & Majority Vote \\
\midrule\noalign{}
\endhead
\bottomrule\noalign{}
\endlastfoot
0 & 0.78 & 0.72 \\
1000 & 0.78 & 0.70 \\
2000 & 0.78 & 0.69 \\
3000 & 0.78 & 0.67 \\
4000 & 0.78 & 0.68 \\
5000 & 0.78 & 0.69 \\
\end{longtable}
}

\textit{[Image omitted: images/758d5524cceff8a39088dd70302986f398c9e4037b775b9b0a65d85418029d63.jpg]}

line

{\def\LTcaptype{none} % do not increment counter
\begin{longtable}[]{@{}lll@{}}
\toprule\noalign{}
Num of Adversarial Labels & GLAD & Majority Vote \\
\midrule\noalign{}
\endhead
\bottomrule\noalign{}
\endlastfoot
0 & 0.78 & 0.72 \\
200 & 0.80 & 0.65 \\
400 & 0.85 & 0.58 \\
600 & 0.90 & 0.50 \\
800 & 0.98 & 0.45 \\
\end{longtable}
}

Figure 5: Accuracy (percent correct) of inferred Duchenne/Non-Duchenne
labels using either GLAD or Majority Vote under (left) noisy labelers or
(right) adversarial labelers. As the number of noise/adversarial labels
increases, the performance of labels inferred using Majority Vote
decreases. GLAD, in contrast, is robust to these conditions.

Results: Using just the raw labels obtained from the Mechanical Turk,
the labels inferred using GLAD matched the ground-truth labels on
78.12\% of the images, whereas labels inferred using Majority Vote were
only 71.88\% accurate. Hence, GLAD resulted in about a 6\% performance
gain.

Simulated Noisy and Adversarial Labelers: We also simulated noisy and
adversarial labeler conditions. It is to be expected, for example, that
in some cases labelers may just try to complete the task in a minimum
amount of time disregarding accuracy. In other cases labelers may
misunderstand the instructions, or may be adversarial, thus producing
labels that tend to be opposite to the true labels. Robustness to such
noisy and adversarial labelers is important, especially as the
popularity of Web-based labeling tools increases, and the quality of
labelers becomes more diverse. To investigate the robustness of the
proposed approaches we generated data from virtual ``labelers'' whose
labels were completely uninformative, i.e., uniformly random. We also
added artificial ``adversarial'' labelers whose labels tended to be the
opposite of the true label for each image.

The number of noisy labels was varied from 0 to 5000 (in increments of
500), and the number of adversarial labels was varied from 0 to 750 (in
increments of 250). For each setting, label inference accuracy was
computed for both GLAD and the Majority Vote method. As shown in Figure
5, the accuracy of GLAD-based label inference is much less affected from
labeling noise than is Majority Vote. When adversarial labels are
introduced, GLAD automatically inferred that some labelers were
purposely giving the opposite label and automatically flipped their
labels. The Majority Vote heuristic, in contrast, has no mechanism to
recover from this condition, and the accuracy falls steeply.

\section{7 Related Work}\label{related-work}

To our knowledge GLAD is the first model in the literature to
simultaneously estimate the true label, item difficulty, and coder
expertise in an unsupervised and efficient manner.

Our work is related to the literature on standardized tests,
particularly the Item Response Theory (IRT) community (e.g., Rasch
{[}10{]}, Birnbaum {[}3{]}). The GLAD model we propose in this paper can
be seen as an unsupervised version of previous IRT models for the case
in which the correct answers (i.e., labels) are unknown.

Snow, et al {[}14{]} used a probabilistic model similar to Naive Bayes
to show that by averaging multiple naive labelers (\textless= 10) one
can obtain labels as accurate as a few expert labelers. Two key
differences between their model and GLAD are that: (1) they assume a
significant proportion of images have been pre-labeled with ground truth
values, and (2) all the images have equal difficulty. As we show in this
paper, modeling image difficulty may be very important in some cases.
Sheng, et al {[}12{]} examine how to identify which images of an image
dataset to label again in order to reduce uncertainty in the posterior
probabilities of latent class labels.

Dawid and Skene {[}5{]} developed a method to handle polytomous latent
class variables. In their case the notion of ``ability'' is handled
using full confusion matrices for each labeler. Smyth, et al {[}13{]}
used a similar approach to combine labels from multiple experts for
items with homogeneous levels of difficulty. Batchelder and Romney
{[}2{]} infer test answers and test-takers' abilities simultaneously,
but do not estimate item difficulties and do not admit adversarial
labelers.

Other approaches employ a Bayesian model of the labeling process that
considers both variability in labeler accuracies as well as item
difficulty (e.g.~{[}8, 7, 11{]}). However, inference in these models is
based on MCMC which is likely to suffer from high computational expense,
and the need to wait (arbitrarily long) for parameters to ``burn in''
during sampling.

\section{8 Summary and Further
Research}\label{summary-and-further-research}

An important bottleneck facing the machine learning community is the
need for very large datasets with hand-labeled data. Datasets whose
scale was unthinkable a few years ago are becoming commonplace today.
The Internet makes it possible for people around the world to cooperate
on the labeling of these datasets. However, this makes it unrealistic
for individual researchers to obtain the ground truth of each label with
absolute certainty. Algorithms are needed to automatically estimate the
reliability of ad-hoc anonymous labelers, the difficulty of the
different items in the dataset, and the probability of the true labels
given the currently available data.

We proposed one such system, GLAD, based on standard probabilistic
inference on a model of the labeling process. The approach can handle
the millions of parameters (one difficulty parameter per image, and one
expertise parameter per labeler) needed to process large datasets, at
little computational cost. The model can be used seamlessly to combine
labels from both human labelers and automatic classifiers. Experiments
show that GLAD can recover the true data labels more accurately than the
Majority Vote heuristic, and that it is highly robust to both noisy and
adversarial labelers.

Active Sampling: One advantage of probabilistic models is that they lend
themselves to implementing active methods (e.g., Infomax {[}4{]}) for
selecting which images should be re-labeled next. We are currently
pursuing the development of control policies for optimally choosing
whether to obtain more labels for a particular item -- so that the
inferred Z label for that item becomes more certain -- versus obtaining
more labels from a particular labeler -- so that his/her accuracy
\(\alpha\) may be better estimated, and all the images that he/she
labeled can have their posterior probability estimates of Z improved.

A software implementation of GLAD is available at
http://mplab.ucsd.edu/\textasciitilde jake.

\section{References}\label{references}

{[}1{]} Amazon. Mechanical turk. http://www.mturk.com.\\
{[}2{]} W. H. Batchelder and A. K. Romney. Test theory without an answer
key. Psychometrika, 53(1):71--92, 1988.

{[}3{]} A. Birnbaum. Some latent trait models and their use in inferring
an examinee's ability. Statistical theories of mental test scores,
1968.\\
{[}4{]} N. Butko and J. Movellan. I-POMDP: An infomax model of eye
movement. In Proceedings of the International Conference on Development
and Learning, 2008.\\
{[}5{]} A. Dawid and A. Skene. Maximum likelihood estimation of observer
error-rates using the em algorithm. Applied Statistics, 28(1):20--28,
1979.\\
{[}6{]} I. Gauthier and M. Tarr. Becoming a ``greeble'' expert:
Exploring mechanisms for face recognition. Vision Research, 37(12),
1997.\\
{[}7{]} V. Johnson. On bayesian analysis of multi-rater ordinal data: An
application to automated essay grading. Journal of the American
Statistical Association, 91:42--51, 1996.\\
{[}8{]} G. Karabatsos and W. H. Batchelder. Markov chain estimation for
test theory without an answer key. Psychometrika, 68(3):373--389,
2003.\\
{[}9{]} Omron. OKAO vision brochure, July 2008.\\
{[}10{]} G. Rasch. Probabilistic Models for Some Intelligence and
Attainment Tests. Denmark, 1960.\\
{[}11{]} S. Rogers, M. Girolami, and T. Polajnar. Semi-parametric
analysis of multi-rater data. Statistics and Computing, 2009.\\
{[}12{]} V. Sheng, F. Provost, and P. Ipeirotis. Get another label?
improving data quality and data mining using multiple noisy labelers. In
Knowledge Discovery and Data Mining, 2008.\\
{[}13{]} P. Smyth, U. Fayyad, M. Burl, P. Perona, and P. Baldi.
Inferring ground truth from subjective labelling of venus images. In
Advances of Neural Information Processing Systems, 1994.\\
{[}14{]} R. Snow, B. O'Connor, D. Jurafsky, and A. Y. Ng. Cheap and fast
- but is it good? evaluating non-expert annotations for natural language
tasks. In Proceedings of the 2008 Conference on Empirical Methods on
Natural Language Processing, 2008.\\
{[}15{]} S. Steinbach, V. Rabaud, and S. Belongie. Soylent grid: it's
made of people! In International Conference on Computer Vision, 2007.\\
{[}16{]} D. Turnbull, R. Liu, L. Barrington, and G. Lanckriet. A
Game-based Approach for Collecting Semantic Annotations of Music. In 8th
International Conference on Music Information Retrieval (ISMIR), 2007.\\
{[}17{]} L. von Ahn and L. Dabbish. Labeling Images with A Computer
Game. In Proceedings of the SIGCHI conference on Human factors in
computing systems, pages 319--326. ACM Press New York, NY, USA, 2004.\\
{[}18{]} L. von Ahn, B. Maurer, C. McMillen, D. Abraham, and M. Blum.
reCAPTCHA: Human-Based Character Recognition via Web Security Measures.
Science, 321(5895):1465, 2008.

\end{document}
